% ════════════════════════════════════════════════════════════════
%  styles.tex —— 自定义样式集中管理（由 common/preamble.tex 在导言区 \input）
%  公式编号 / 代码框 / 推导框 / 自定义定理类单元 / 新单元引用名
% ════════════════════════════════════════════════════════════════

% ---- 公式编号：ElegantBook 默认全书连续，改为「章.序」(式 3.1) ----
\numberwithin{equation}{chapter}
\numberwithin{figure}{chapter}   % 图按「章.序」编号（图 2.1）
\numberwithin{table}{chapter}    % 表按「章.序」编号（表 2.1）；小结表加 \caption 后随之编号

% ---- 允许 >10 列矩阵（amsmath 默认上限 10）：如 DLT 12 列设计矩阵 ----
\setcounter{MaxMatrixCols}{20}

% ---- 章节编号/目录深度：只编到 section，避免 \paragraph 出现 2.2.0.0.1 式深层号 ----
\setcounter{secnumdepth}{2}
\setcounter{tocdepth}{1}

% ---- 章首页（含目录/前言等走 \chapter*→plain 样式的页）显示页码 ----
%      ElegantBook 把 plain 重定义为 \fancyhf{} 全清空 → 每章首页、目录首页都无页码。
%      此处覆盖：保留与正文一致的「页脚居中页码」(structurecolor\small\thepage)，仅清页眉与页眉线。
\fancypagestyle{plain}{%
  \fancyhf{}%
  \fancyfoot[C]{\color{structurecolor}\small\thepage}%
  \renewcommand{\headrulewidth}{0pt}%
  \renewcommand{\headrule}{}%
}

% ---- \paragraph 小标题「标题前距」收紧：book.cls 默认 3.25ex(≈1.4行) → 1.2ex(≈半行) ----
%      正文大量用 \paragraph{...} 作段首粗体小标题，默认前距过松（视觉上"整空行"）。
%      仅缩 beforeskip，保留「粗体随行小标题 + -1em 后距(run-in)」原样；
%      elegantbook/ctexbook 均未重定义 \paragraph，故直接照 book.cls 原式重定义、只换前距。
\makeatletter
\renewcommand\paragraph{\@startsection{paragraph}{4}{\z@}%
  {1.2ex \@plus.5ex \@minus.2ex}{-1em}{\normalfont\normalsize\bfseries}}
\makeatother

% ---- 代码框需要 tcolorbox 的 listings 库（类已载 tcolorbox + listings 宏包）----
\tcbuselibrary{listings}

% ---- listings 代码风格（配主题色 main/second）----
\lstdefinestyle{notecode}{
  basicstyle=\ttfamily\footnotesize,
  keywordstyle=\color{main}\bfseries,
  commentstyle=\color{gray!70}\itshape,
  stringstyle=\color{second},
  numbers=left, numberstyle=\tiny\color{gray}, numbersep=8pt,
  breaklines=true, showstringspaces=false, tabsize=2, columns=flexible,
}

% ---- 纯文本 / 伪代码语言（listings 无内置 text）：空定义，无关键字高亮 ----
%      用于规划/控制章的算法伪代码框 \begin{codebox}[text]{...}
\lstdefinelanguage{text}{}

% ---- ① 代码框（区分“代码”内容）：带标题彩框 ----
%      用法：\begin{codebox}[C++]{标题} ... \end{codebox}
%      可选 #1 = 语言名（listings 语言，如 C++/Python/bash/Matlab）：驱动语法高亮，
%      同时作为右上角语言角标（main 色，与左上标题同色；#1 留空则角标自动消失）。
%      几何：boxsep/top/bottom=3pt + middle=0pt —— 去掉 tcblisting 为下段预留的空带，
%      绿框紧贴代码、内层蓝线（继承 elegantbook 全局 lstset frame）上下居中。
\newtcblisting{codebox}[2][]{%
  listing only, listing engine=listings, listing style=notecode,
  listing options={language=#1},
  enhanced, breakable,
  colback=gray!6, colframe=main!75!black, coltitle=white,
  fonttitle=\bfseries\small, title={#2},
  attach boxed title to top left={xshift=8pt, yshift=-2pt},
  boxed title style={colback=main!75!black},
  left=10pt, boxsep=3pt, top=3pt, bottom=3pt, middle=0pt,
  % 左上标题用 tcolorbox 原生 boxed title（自动预留上方空间，不挤压上文，无连锁位移）。
  % 右上角语言角标(#1)：与原生标题块同字号(\bfseries\small)/同色/同高度
  % （minimum height=\tcboxedtitleheight 即取原生标题块高度），锚到其垂直中心，
  % 仅左右对齐相反；#1 空则不画。
  overlay unbroken and first={%
    \if\relax\detokenize{#1}\relax\else
      \node[anchor=east, font=\bfseries\small, text=white, fill=main!75!black,
        rounded corners=2pt, inner xsep=6pt, inner ysep=0pt,
        minimum height=\tcboxedtitleheight, xshift=-8pt]
        at (title.center -| frame.east) {#1};%
    \fi},
}

% ---- ② 推导框（区分“推导过程”内容）：可跨页 ----
%      用法：\begin{derivation}[可选标题] ... \end{derivation}
\newtcolorbox{derivation}[1][推导]{%
  enhanced, breakable,
  colback=third!4, colframe=third!70!black, coltitle=white,
  fonttitle=\bfseries\small, title={#1},
  attach boxed title to top left={xshift=8pt, yshift=-3pt},
  boxed title style={colback=third!70!black},
  boxrule=0.5pt,
}

% ---- ③ 自定义定理类单元（官方接口 {名}{显示名}{风格}{引用前缀}）----
%      风格：defstyle(绿/main)  thmstyle(橙/second)  prostyle(蓝/third)
\elegantnewtheorem{algo}{算法}{prostyle}{alg}
\elegantnewtheorem{insight}{洞见}{defstyle}{ins}
\elegantnewtheorem{paper}{论文}{thmstyle}{paper}   % 橙 论文精读盒（可 \cref）

% ---- ⑤ 陷阱盒（概念误区 / 思维陷阱）：红框警示，无计数器（不参与 \cref）----
%      用法：\begin{pitfall}[可选标题] ... \end{pitfall}
\newtcolorbox{pitfall}[1][陷阱]{%
  enhanced, breakable,
  colback=red!3, colframe=red!55!black, coltitle=white,
  fonttitle=\bfseries\small, title={#1},
  attach boxed title to top left={xshift=8pt, yshift=-3pt},
  boxed title style={colback=red!55!black},
  boxrule=0.5pt,
}

% ---- ⑥ 练习 / 前置自测盒：橙框，无计数器 ----
%      用法：\begin{practice}[可选标题] ... \end{practice}
\newtcolorbox{practice}[1][练习]{%
  enhanced, breakable,
  colback=second!4, colframe=second!60!black, coltitle=white,
  fonttitle=\bfseries\small, title={#1},
  attach boxed title to top left={xshift=8pt, yshift=-3pt},
  boxed title style={colback=second!60!black},
}

% ---- ④ cleveref 中文引用名 ----
%      ElegantBook 定理盒基于 tcolorbox，计数器名为 tcb@cnt@<环境>，须按此绑定，
%      否则 \cref 报 “cref reference format for label type 'tcb@cnt@...' undefined”。
\makeatletter
\crefname{tcb@cnt@theorem}{定理}{定理}
\crefname{tcb@cnt@definition}{定义}{定义}
\crefname{tcb@cnt@proposition}{命题}{命题}
\crefname{tcb@cnt@lemma}{引理}{引理}
\crefname{tcb@cnt@corollary}{推论}{推论}
\crefname{tcb@cnt@axiom}{公理}{公理}
\crefname{exam}{例}{例}        % ElegantBook example 用 amsthm 计数器 exam（非 tcb@cnt@example），否则 \cref 显示 ??
\crefname{tcb@cnt@exercise}{练习}{练习}
\crefname{tcb@cnt@problem}{习题}{习题}
\crefname{tcb@cnt@algo}{算法}{算法}
\crefname{tcb@cnt@insight}{洞见}{洞见}
\crefname{tcb@cnt@paper}{论文}{论文}
\crefname{figure}{图}{图}        % cleveref 默认 figure 走英文 'figure'，此处本地化为中文「图」
\Crefname{figure}{图}{图}
\crefname{table}{表}{表}          % 同理本地化 table（全书级一致）
\Crefname{table}{表}{表}
\makeatother
